\chapter{适应性}

训练是以身体被施加某种刺激，并从该刺激中恢复，然后适应这个刺激的过程为基础的，生命的进程因此得以在存在这种刺激的环境中延续。这是一个极为基本的生物学概念，对刺激的适应能力是定义生命的重要标准之一。理解这一点对想要富有成效地训练而不只是锻炼的教练和运动员来说至关重要。这样的理解始于汉斯·塞利（Hans Selye）博士。

在1936年7月4日，《自然》（\textit{Nature}）杂志发表了一篇名为《由多种有害介质引发的综合征》（\textit{A Syndrome Produced by Diverse Nocuous Agents}）的论文。该论文的前提是，一个生命体暴露在外部的压力源（刺激）中，之后会表现出一系列特定的短期反应和长期适应。在我们的讨论中，这个刺激就是举重。

\section{一般适应综合征}

塞利把锻炼看作一种“有害”或者有毒的压力源，在负荷过大或频率过高时甚至可能致死。他的理论来自于对应激动物的观察和对应激细胞的显微学观察。他的研究是在对人体代谢和骨骼肌收缩机制的细节完全不了解的情况下进行的，甚至在论文发表时这些还没有搞清楚。尽管支持自身论点的信息相对匮乏，但他的推论非常合理。现如今，我们对生理学机制的理解更为完整，从而可以更好地理解和应用塞利的理论。我们对急性期反应和后续的适应性反应的理解，既包含非常明确的时间过程，又加上了对应激后细胞事件的现代见解，这为塞利的先见之明增加了不少分量。

塞利提出，重复暴露在非致死程度的压力源下会使身体在随后再次暴露于同种刺激时产生耐受性，因为对特定刺激的适应具有对该刺激的针对性。这就为特异性这个概念提供了支持——训练刺激必须与训练希望达到的运动表现相关才能产生可以提高特定运动表现的适应状态。这个理论认为，身体会经历三个可能的阶段，前两个阶段是为了生存，第三个阶段则意味着身体无法承受或继续适应外界的刺激。

\textbf{第1阶段：警觉或惊醒（Alarm or Shock）}\quad 警觉阶段是刺激出现之初身体的即时反应，期间会发生很多事情。塞利指出，第1阶段的一个重要特点就是“肌张力”（Muscular Tone）的快速丧失可持续48小时左右。现在我们知道，这个阶段发生的其他过程还包括炎症反应和急性期反应，并且这些效应都是针对导致其产生的特定刺激的（也就是说，手上被烫伤不会导致脸上起泡，10英里长跑不会造成手部肌肉酸痛）。后两种反应的主要结果之一是，在压力源消退前，基本的细胞活动会受到抑制，以稳定细胞结构及其代谢水平。这是一个争取生存的过程，同样也可以作为有效训练刺激的标志。这个阶段也可能伴随有轻微的肌肉和骨骼的不适，这意味着体内平衡的打破和训练后肌肉内刺激其结构和功能改变的事件正在发生。训练者在这个阶段可能不会感受到酸痛或疼痛，他更可能的感受是“僵硬”或者“疲劳”。无论主观感受如何，这个阶段会出现短暂的运动表现下滑，虽然这种下滑在典型的5磅（2.3千克）增幅的杠铃训练体系中很难察觉。运动表现的下滑在基于技巧和爆发力的训练中会更为明显，在最大力量的训练中则不易被觉察。

塞利没有料到，他的理论会成为健康个体制订训练计划的核心依据。如果他当时意识到了他的理论对运动员训练的重要性，那么基于机体当前的适应状态，第1阶段可能会得到更加深入的阐述。对于初级训练者，由于训练还未能提高其力量或训练承受度，因此初级训练者应使用比高级训练者更小的重量来打破体内平衡。随着训练者水平的提高（从初级到中级再到高级），刺激的大小或持续时间也需要随之增长才能促成第1阶段的出现。



\textbf{第2阶段：适应或抵抗（Adaption or Resistance）}\quad 在第2阶段，身体通过调节基因活性、改变激素产量、增加结构性和代谢性蛋白质的合成来响应训练刺激，这些过程的累积效应就是我们所说的“恢复”。实际上，身体会试图提高自己承受反复刺激的能力以确保生存。刺激效应会产生一系列针对刺激的反应，并从该刺激中得到特定的恢复。在训练环境中，当这个过程发生时，运动表现就会提高。塞利总结道：适应阶段通常出现在身体接受刺激后2天左右，并且如果周期性地施加同样的刺激，那么在4周或更短的时间内身体就会产生完全的适应状态。

我们现在明白，训练承受能力决定了一个人接近其最终身体潜力的程度，而适应状态的发生会根据每个人当前的训练承受水平的不同变化。一些距离最终身体潜力还很远的人（初级训练者）能够很快产生适应状态，这个过程会从受到刺激时开始，并在24～72小时完成。对这类人来说，不需要太大的刺激就足以打破其体内平衡，并且即使状况没有达到最优，也很容易在指定时间内得到恢复。在另一种极端情况中，一个高级训练者可能需要1～3个月，甚至更长时间，来适应一次足够大的、业已超过其高度发达的训练承受能力的、足以打破其体内平衡并促使进一步的适应状态出现的刺激。

很重要的一点是，这个系统处于动态的变化中，不同层面上有无数的事情可以作为压力源刺激适应过程的发生。适应状态也会出现在先前已经适应的刺激因素被移除的时候，所以，根据引发适应状态的事件不同，适应状态可以出现在各个“方向”。在这种情况下，刺激可以是任何能够引发适应状态的事件。

另一个很重要的点是，身体有理由也有能力根据外界情况来改变自身的准备状态：如果没有外界刺激的存在，持续保持高度的准备状态代价极其高昂。我们总是无法达到我们的遗传潜力所赋予的强壮和快速，是因为周围的环境条件并不总是需要那样的状态，保持这样的状态需要调动大量宝贵的代谢资源，而这些资源在较低的准备程度下经济使用的话可以持续更长的时间。同理，在易变的环境中无法适应的基因就会被淘汰。准备程度的提高需要对刺激做出反应，从而使机体处于准备程度高于先前的状态中。形成高于先前基线的适应能力在生命的早期就已进化出来。

\textbf{第3阶段：力竭（Exhaustion）}\quad 不论是强度、持续时间还是频率，如果作用于身体的刺激过大，身体就无法充分地适应，力竭就会出现。塞利提出，持续1～3个月的过度刺激可能会导致个体死亡。在1936年，用这种武断的估计来阐述理论没有什么问题，但基于我们对运用这项原则的现代理解，以及跑3小时马拉松也存在猝死案例的事实，我们认为，这个论点并不完全适用于训练。如果我们认同不适量的训练刺激会导致力竭，那么对不同的训练进步幅度来说这个刺激的大小也应不同。在现实中，主要是中高级训练者需要考虑这一点——初级训练者往往缺乏足够的力量和耐力，其训练强度和持续时间无法产生这种程度的刺激（不过没有经验的教练胡乱指导新手训练仍可能会导致力竭的情况出现）——也就是应当避免持续时间过久、强度过大的训练。底线是，没有人想要达到第3阶段，也就是所谓的“过度训练”阶段。

塞利的理论在训练中的实际应用如图\ref{fig:2-1 塞利的理论在训练中的实际应用}所示。在一般情况下，渐进性训练需要在身体产生了明显的适应状态后尽快增加训练负荷。一直使用最初的、身体已经适应的负荷无法打破体内平衡，因为适应状态已经存在了，重复使用相同的训练负荷无法带来持续的进步。如果你的目标是提高运动表现的话，使用一成不变的训练负荷是一种毫无效果（但却常见）的执教和训练方式。

\begin{figure}[htbp] % htbp 控制图片排版位置优先顺序
    \centering
    \includegraphics[width=0.8\textwidth]{image_001.jpg}
    \caption{基于塞利的理论参数，身体在接受训练刺激后存在三种可能的发展路径。刺激太小无法打破体内平衡，也就不会产生什么改变；刺激过大虽然会打破体内平衡，但却超出了身体的适应能力，从而导致表现下滑；只有适量的刺激既能够打破体内平衡，同时又不会超出机体的适应能力，才能促成进步}
    \label{fig:2-1 塞利的理论在训练中的实际应用}
\end{figure}

能够在塞利的理论基础上理解锻炼和训练的区别十分重要。刺激—恢复—适应的过程是每个人都可用于自己的体育训练中来提高其运动表现的符合逻辑的过程。训练者只须持续地对身体施加与他希望提高的运动表现相关的、可以引发适应状态的刺激。根据定义，进步需要变化。因此，不断变化的刺激是有效训练的天然要素。

这就意味着，在最基本的层面上，长跑训练和举重训练是完全不同的。这两项运动所需要的运动表现可以说是两个极端，为了获得达到冠军表现所需的适应状态，训练施加的刺激必须体现出这种不同。对这两种运动项目来说，有效的训练都要求施加与各自的项目对应的、渐进的刺激。举重选手无法通过长跑提高其举重能力，同样，长跑运动员强调力量训练也无法缩短长跑时间。相比耐力，力量是一种更为基础的适应，如果根据比赛日程合理安排力量训练，长跑运动员也会从中受益，但竞技型的举重运动员则没有必要在备赛计划中加入长跑训练。

如果一名举重选手想要变得更加强壮，那他必须做好计划；如果一名长跑运动员想要跑得更快，他也要做好计划。举重运动员通过计划举起更大重量的方式来变强壮，长跑运动员则通过计划在更短时间内跑更多里程的方式来提高速度。产生这两种适应的一般过程相当简单，尤其是对那些刚开始训练举重或长跑的训练者来说，如果采用一种可持续进步的方式，他们提高的空间很大，他们可以举起更大的重量、跑得更快。必须施加刺激，并从该刺激中恢复，由此产生的净效应就是可以提高运动表现的适应状态。但是人们往往会忽视，取得进步的过程是由长期逐渐增加的刺激—恢复—适应效应累积而来的，单次的训练本身并不是关键。随着时间的推移，累积的训练效应才会产生适应状态。如果没有控制好这一点，训练就会缺乏效率；如果完全不加控制，那我们所做的就不是训练，只是锻炼罢了。

在不同情况下，对于不同的人，刺激—恢复—适应循环的作用方式也不尽相同。初级训练者的终极身体潜力还完全没有开发，对他们来说，从一次力量训练中恢复并产生适应状态后，他就能够变得更加强壮。从这个点开始，只要没有因为不再施加刺激使得先前的适应状态消失，后续施加的刺激就能使这个过程再次发生，并产生进一步的适应状态。对已经经历过这个过程，累积了多年的适应状态，业已变得很强壮的运动员来说，这个过程的发生则要慢得多，并最终达到一种几乎无法继续获得适应状态的地步（见图\ref{fig:1-1 力量表现提高与训练复杂度相对于时间的一般关系}）。从初级训练者的第一次训练到高级运动员为了增加1千克总成绩而倾尽全力，这个过程存在一个连续的范围，每个训练者都能在其中找到自己的位置，他们所处的位置决定了他们对刺激—恢复—适应过程的响应水平。

对初级训练者来说，单次训练就能打破肌肉和身体系统内的生理平衡。如果一次力量训练就可以打破体内平衡，根据打破的程度，可以详细推得一组结果。图\ref{fig:2-1 塞利的理论在训练中的实际应用}描述了这个模型。

对初级训练者来说，一次适度的对生理平衡的打破会导致短暂而轻微的运动表现下滑。下滑幅度很小是因为初级训练者的水平本身非常低，这个级别的小幅下滑很难测量到。这种下滑会在训练后立刻出现，表现为塞利模型的第1阶段。训练后的数小时至几天内，运动表现会恢复到正常水平，然后超越施加刺激之前的水平。这被称作超量恢复（Supercompensation），一种身体已经准备好了承受比其正在准备的以及已经成功经历的更强刺激状态的过程。此时，训练者已经成功完成了塞利模型中的第2阶段，适应了初始的训练负荷（见图\ref{fig:2-2 超量恢复曲线}曲线A）。

\begin{figure}[htbp] % htbp 控制图片排版位置优先顺序
    \centering
    \includegraphics[width=0.8\textwidth]{image_002.jpg}
    \caption{训练者从初级阶段到高级阶段的成长过程需要较长的刺激—恢复—适应周期。（A）初级训练者能够通过一次训练获得足够的刺激，然后在72小时内超量恢复到基线之上。（B）在1周的微周期内，中级训练者需要几次训练和更长的恢复时间才能引发超量恢复。（C）高级训练者需要多次训练累计的刺激以及1个月甚至更长时间的恢复才能取得进步}
    \label{fig:2-2 超量恢复曲线}
\end{figure}

你要明白，训练者并不是在训练过程中变强壮的，而是在训练后的恢复阶段变强壮的。那么，既然他现在变强了一些，符合逻辑的做法就是在下次训练时增加训练负荷，也就是采用简单的渐进性超负荷方式——有规律地逐渐增加最大训练重量。再次使用同样的训练负荷不会产生任何提高，因为训练者对这个负荷已经适应了。这时，增加一点训练负荷就能再次让训练者经历塞利第1阶段和第2阶段，使超负荷—适应循环在一个稍高一点的水平上重复。如果每次训练增加的负荷相同，那么我们就将其称为\textbf{线性增长计划}。

这样的训练安排能持续使用几个月，直至训练遇到平台期。这时，很可能需要以周为单位安排一系列的训练，每周专门安排2～3次训练以产生累积效应，形成较长的训练—恢复周期，才足以使训练者度过塞利模型中的前两个阶段，这代表着中级训练者的响应情况（见图\ref{fig:2-2 超量恢复曲线}曲线B）。根据训练者的目标的不同，其训练生涯的中级阶段可能会很长，也可能只有数年。

随着身体能够更好地对抗重量发力，接受刺激并从中恢复的能力也会越来越强。运动表现和恢复能力会随着渐进式训练的推进而增长，最终需要几周才足以打破体内平衡以刺激适应状态的产生，同时需要很长的时间从中恢复并产生超量恢复。高级训练者甚至需要1个月的时间才能度过第1阶段和第2阶段（见图\ref{fig:2-2 超量恢复曲线}曲线C）。

运动员的训练水平越高，理解刺激—恢复—适应模型就越重要，理解该模型平衡构成人体适应状态的两种对立要素的方法就越重要：（1）训练负荷必须足以打破体内平衡才能引发适应状态的产生，同时训练负荷不能过高，变成一种无法承受的刺激；（2）必须保证充分的恢复才能促使适应状态出现（见图\ref{fig:2-3 训练负荷与训练承受度的关系}）。高级运动员行走在刀刃上。对初级训练者来说，他们的面前是一条康庄大道，基本看不到什么刃口，行走起来很容易；对中级训练者来说，这条路要险恶一些，需要使用更复杂的方法；对高级训练者来说，想要在这锋利的刀刃上保持平衡而不受伤，他们需要非常仔细地控制训练计划的每一个变量。

\begin{figure}[htbp] % htbp 控制图片排版位置优先顺序
    \centering
    \includegraphics[width=0.5\textwidth]{image_003.jpg}
    \caption{训练负荷与训练承受度的关系。不论处于哪个训练阶段，都存在一个训练负荷上限（以小箭头指示），超出这个界限就会导致过度训练，造成训练承受度和运动表现下降。注意，（1）训练承受度在训练生涯中会显著提高，承受度的增加源于训练负荷的渐进式增加。（2）随着训练生涯的推进，过度训练导致的承受度下降和运动表现的退步会表现得更快速、更突然（注意图中的曲线，在训练承受度达到峰值后，级别越高的训练者，其承受度的降幅越大）。虽然过度训练对每个阶段来说都是个问题，但在高级训练阶段防止过度训练尤为关键，因为在该阶段，一旦训练强度超过了训练者的承受度，其运动表现会急剧下降。相对而言，在初级训练阶段和中级训练阶段判断是否过度训练本身成了一个问题，因为过度训练导致的运动表现下滑的速率要慢得多，很容易被忽视或错误解读}
    \label{fig:2-3 训练负荷与训练承受度的关系}
\end{figure}

\section{解读过度训练}

\textbf{超负荷}（Overload）这个概念是解读在力量训练中取得进步的关键。超负荷代表着打破体内平衡和引发适应状态所需的训练负荷。想要取得进步，必须使生理系统受到扰动，在力量训练中，这个扰动就是更大的重量、更大的训练量（由组数和重复次数确定），或者对中级训练者和高级训练者来说，可以是更短的组间休息时间。通过训练，将超负荷施加在生理系统上，以特定的方式打破体内平衡，这个过程被称为\textbf{超负荷活动}（Overload Event）。对初级训练者来说，每次训练都构成一次超负荷活动。对中高级训练者来说，在1周或更长的训练时间内，所有加码的要素共同构成了超负荷活动。

但如果身体无法从一次超负荷活动中得到恢复，超负荷就无法产生进步。缺乏充分恢复的超负荷活动会引起过度训练。\textbf{微周期}（Microcycle）这个术语在传统上被定义为1周的训练，更好的定义方式是超负荷活动和从超负荷活动中恢复所需要的时间。这段时间会因训练者水平的不同而变化。对初级训练者来说，一个微周期就是两次训练之间的时间间隔。训练者水平越高，微周期就会越长，直到这个词对高级训练者失去意义，因为高级训练者这个过程所需要的时间一般被称为\textbf{小周期}（Mesocycle）。对我们来说，这些术语不够简明，缺乏实用性。

过度训练是任何训练计划的灾难。当训练刺激超过身体适应能力的上限时，训练者不仅会停滞不前，甚至可能退步。用我们讲过的术语描述就是，过度训练的运动员进入了塞利第3阶段。训练量或训练强度与身体恢复之间出现了失衡，从而导致训练者无法从刺激中恢复，超量恢复也不会发生。疲劳的影响很大，以至于恢复过程受到影响，甚至无法得到恢复，导致持续疲劳，甚至疲劳程度加剧。运动表现会一直保持初始超负荷后的低迷状态，并随着训练的继续进一步下降。最终的结果就是无法训练，无法表现出原有的水平。

在传统的运动科学文献中，运动刺激有三种可能的版本：（1）疲劳（Fatigue）；（2）超量训练（Overreach）；（3）过度训练（Overtrain）。每一种都与运动表现的下降相关，但只有一种是真正的问题。

\subsection{疲劳}

在生理学中，疲劳通常被定义为肌肉生成力量能力的下降。疲劳可以仅仅是身体活动后产生的简单和短暂的“疲惫感”——由进入塞利第1阶段所必需的刺激形成的训练要素。在最基本的刺激—恢复过程中，初级阶段的疲劳应当经过48～72小时得到恢复。中级训练者的疲劳要素会在完成一个完整的训练周后出现。而对高级训练者来说，完全恢复至超量恢复状态可能会在疲劳仍存在的状况下出现，而且可能需要1个月甚至更久。对中高级训练者来说，每次训练开始时完全感觉不到疲劳并不现实，也不需要这样。如果一个训练者持续感受不到疲劳，那说明训练计划不够严格，无法引发体内平衡的扰动和适应状态的产生。

\subsection{超量训练}

超量训练是指一系列训练的累积效应使训练者出现了短期的运动表现下滑、感觉疲惫、心情郁闷、身体疼痛、睡眠质量不佳和其他各种需要2周以上才能恢复的情况。此时体内平衡的紊乱也会导致短期的睾酮水平降低、皮质醇水平升高等激素水平的变化。其实，这些也是产生积极的杠铃训练系统效应的部分要素。但这个定义有个明显的问题，即超量训练与过度训练的差别仅仅在于，前者可以通过大约2周的减量训练或休息得到恢复，而过度训练的恢复则需要更长时间——这种区分方式显然太过随意。

并且，对超量训练如此定义没有考虑到个体训练进步水平的差异以及由此带来的恢复能力的差异，对这些个体差异的忽视也是传统运动科学文献的典型问题。只要有一个设计得当的新手计划，初级训练者是不会经历超量训练的，因为他们只需要48～72个小时就能实现超量恢复。除非使用了远超建议的负荷，初级训练者是不会也不可能超量训练的，因为初级训练者的标志性特点就是能够从负荷渐增的训练中迅速恢复，从而产生渐进的、稳定的进步，理解这一点非常重要。即便是在初级训练阶段后期使用简单线性计划进入平台期时，额外安排一次短期的减量训练通常也足以使运动员的恢复能力回归正常。中级训练者能够在1周的训练安排中恢复，因为中级训练者的典型特征就是能够响应短期累积的训练负荷。经历了很长时间的体内平衡的累积扰动之后，高级训练者需要4周或更长时间才能恢复并产生超量恢复，这比超量训练定义的2周时间要长，但这只是高级训练者正常的训练计划时间框架。因此，“超量训练”这个概念对力量训练并没有什么用处。

另外，将“超量训练”定义为一种负面的训练影响实在是莫名其妙。一名运动员必须“超量训练”才能产生足够的刺激来打破体内平衡，必须有意使用超过其已经适应的最大负荷的训练量才能引发超量恢复。抛弃“超量训练”，使用“超负荷”的表述要更贴切、更实用并易于理解，因为“超负荷”描述了所有训练者产生适应状态所需要的负荷和刺激，并且不需要考虑进步程度。所有的训练计划都应当包含超负荷阶段，因为这是实际运用塞利理论所要求的；我们应当把这些阶段理解为适应性的，而不是有害的；并且如果想要得到期望的效果，应当根据训练者的进步水平合理安排超负荷。判断超负荷的幅度并非易事，这需要极强的驾驭能力，对高级训练者来说这一点尤为重要，因为负荷过大或者恢复不足会导致运动表现快速下降甚至过度训练。

\subsection{过度训练}

过度训练是在缺乏充分恢复的情况下使用过高训练量或训练强度，或者兼具二者的累积结果，过度训练会消耗身体从训练刺激中恢复并产生适应状态的能力。判断过度训练的首要指标是，在进行了通常足以恢复的休息之后，运动表现仍然没有提高。尽管人们（美国运动医学会和美国奥林匹克委员会）普遍接受的定义认为，发生过度训练时需要不少于2周的休息才能恢复过来，但这显然与训练者的水平有关，不可能有一个硬性的、统一的标准来判定并避免过度训练的发生。尽管让初级训练者使用过大的训练负荷不可取，会使其立刻失去运动表现能力，但他们仍能很快恢复。虽然没有观察足够的时间，但教练所看到的症状就是过度训练的表现。尽管初级训练者也会出现过度训练，但由于初级训练者缺乏可以对比的训练史并且整体水平很低（如图2-3所示），其运动表现下滑的幅度难以察觉，所以初级训练者是否过度训练并不容易判断。和“超量训练”一样，中级训练者的过度训练适用于美国运动医学会和美国奥林匹克委员会的定义：使中级训练者无法在2周内恢复的训练则视为过度训练。但对高级训练者来说，因为其训练周期更长，这样的定义对他们已不再适用。不过，判断高级训练者是否过度训练相对容易，因为基于他们长期的训练史，其运动表现的下滑程度很容易看出来。

一个对所有训练阶段都适用的过度训练的定义需要找到能够合理量化各个阶段恢复时间的方法。\textbf{当运动表现在一个减重训练周期后仍无法恢复时，过度训练就发生了。}这个周期的长度取决于运动员的训练水平。比如，一个初级训练者每48小时训练1次，然后在一次训练中由于先前训练量过大导致状态非常差（这在热身阶段表现得很明显）。因为身体疼痛，他的动作幅度变小了，移动杠铃的速度明显变慢，增加重量后感觉更为吃力。这时，教练就应当阻止其继续训练，找出问题所在（比如，他在上次训练时趁着教练去另一个屋子指导别人的时候额外做了5个正式组），让他回家休息，直到48小时后开始下一次训练。训练者回归训练后，其热身状态显示他已经恢复了，现在能够完成本该上次完成的正式组训练了。这个过程说明，他之前过度训练，现在已经恢复了。可能因为他是一个初级训练者，不论是正常的超负荷还是过度训练，他都能快速得到恢复（初级训练者的特点），因为内在的机制是相同的。

如果一个致力于4周周期的高级训练者表现下降，低于预期水平，该运动员要么是进入了过度训练的周期，要么是当前的训练周期已经耗尽了他的恢复能力。这种情况下，需要多至4周的削减重量的训练来帮助恢复。无论是初级训练者还是高级训练者，当其被判定为过度训练后，为了尽快恢复体内平衡，都需要重复一个与正常训练周期等长的、大幅削减训练负荷的训练周期。高级训练者使用的训练周期一般较长，花几个月甚至更长时间去发现并纠正计划中的错误，这个代价他们承受不起。

过度训练也是初级训练者与高级运动员之间的显著差别之一。运动员水平越高，过度训练的代价就越大。初级训练者可能会因为错过一次训练或者因为某次训练目标没有达成而感觉不佳，但这只会持续几天，并且除了对下一次训练有些许影响之外不会造成任何后果。中级运动员已经投入了很多，到了选择自己的运动专项进行训练的时候，并在为成为竞技运动员而努力。高级运动员就是为了比赛而训练的人，他们已经为比赛投入了数千小时的训练、无数的金钱和汗水，并甘愿为此付出牺牲。精英级运动员可能有冠军头衔、赞助费、广告代言，每次比赛都可能决定他们退役后的生活水平。随着职业生涯的发展，失败的代价会越来越高，即使暂时的失败也可能带来严重的后果。

我们对过度训练过于重视了吗？根据美国奥林匹克委员会和美国运动医学会“对过度训练的共识”，每天都有10\%～20\%的运动员会过度训练。如果这是真的（当然不太可能，因为大多数运动员的训练根本没有刻苦到过度训练的地步，美国奥林匹克委员会和美国运动运动医学会应当知道这些），这就是个问题了。

有多少教练能够承受，在比赛当天20\%的队员表现低于日常水准？无论何时，如果有这么多运动员存在过度训练，都会严重影响团队的成功，对个人职业生涯来说也是如此。过度训练是运动员没有真正理解并在训练中正确运用刺激—恢复—适应理论的后果，也是没有对运动员的训练水平和进步进行正确评估的后果。

对非初级的训练者来说，过度训练的征兆一旦显现，症状会非常严重：表现明显下滑、睡眠质量差、慢性疼痛增多、情绪异常波动、慢性的心率升高、食欲不振、体重下降以及其他各种生理和心理异常（实际上，这与严重抑郁症的生理症状相同，抑郁症也是由长期的压力积累导致的）。不过，即使使用同样的计划，同样出现了过度训练，所有的训练者也不会表现出同样的症状。再强调一遍，教练的眼睛是判断运动员的表现变化和健康状态必不可少的。一旦确认了过度训练，采取补救措施刻不容缓，因为过度训练出现的时间越长，恢复所需的时间也会越长。很可能获得恢复需要付出的时间会2倍于造成过度训练的时间。时常会听到运动员因过度训练而错过了整个训练年的恐怖故事。因此，你必须全力以赴，识别和处理这种严重的情况。

\section{影响恢复的因素}

对过度训练的讨论往往非常狭隘，基本只讨论训练与恢复的时间比例。尽管它们确实是打破体内平衡和产生力量适应的两个控制因素，但是，恢复毕竟是一个受多因素影响的事件，它不是只受到训练间隔时间的影响。经常会在铁杆儿举重区和健身社区里听人讲“噢，根本不存在过度训练！”首先我很确定，过度训练是存在的，但这些人的态度表明他们意识到了，还有很多其他因素有助于恢复，因此避免过度训练是完全可能的。注意恢复期的饮食和休息的细节对避免过度训练至关重要。除非教练和训练者都理解并努力实现最优恢复，否则任何训练方法都无法产生最佳效果或避免过度训练。

除了训练和休息的时间比例，还有很多影响或促进恢复的因素，其中最重要的就是足量的睡眠和饮食——适量摄入蛋白质、热量、水和微量营养素。问题在于，这些因素都是由训练者——而不是教练——直接控制的。一个优秀的教练会向训练者解释，这些事情为什么重要，并会定期强调它们的重要性，让训练者意识到，出色的运动员会把这些事情当作自身应尽的责任，而平庸的运动员则不会。如果运动员无法将自己该做的事情做好，即使用全宇宙最好的训练计划也会功亏一篑。任何训练计划的成功最终都取决于运动员的责任感。

\subsection{睡眠}

睡眠的重要性显而易见，但在身体需求和刺激增加的阶段，教练和训练者却经常会忽视睡眠的必要性。不论怎样强调睡眠对力量型运动员的意义都不为过——这可能是我们可以控制的最重要的合成代谢因素。虽然适用范围有限，但这一主题的科学文献支持以下观察。

\begin{enumerate}
    \item 恢复期间缺乏充足的睡眠会导致竞争能力、决断能力以及对训练强度的承受度下降。
    \item 缺乏充足的睡眠对情绪状态有负面影响，导致主观疲劳感增强、抑郁症状加重，甚至导致轻微的混乱。
    \item 缺乏充足的睡眠会导致身体响应训练刺激、产生适应状态的生理能力下降。
\end{enumerate}

睡眠期间会发生很多生理变化。其中，激素分泌可能是机体为了恢复做出的最重要的努力。睡眠期间，参与合成代谢（肌肉合成）的激素浓度升高，参与分解代谢（肌肉损耗）的激素浓度下降。一入睡睾酮水平就会开始升高，并在第一次快速眼动（Rapid Eye Movement，简称REM）时达到峰值，且在醒来之前一直保持这个水平。这意味着，睡眠模式一旦被干扰很可能会阻碍睾酮对恢复的促进作用。另一种合成代谢激素——生长激素，在睡眠时期也有其典型的分泌模式。在深度睡眠开始不久后，生长激素浓度开始升高，并在达到峰值后继续维持1.5～3.5小时。生长激素的一个主要功能就是抵消分解代谢激素皮质醇的负面作用。睡眠的中断或睡眠时间的缩短会减少这些重要的合成代谢激素的有益效果。

那么需要睡多久呢？美国军方曾经认为，每晚连续睡4个小时就足以维持生存和基本的战斗能力。但现在，我们意识到，我们需要更多的睡眠时间，建议每晚睡7～8个小时来保持“持续的运行状态”。妈妈会告诉我们，每天晚上睡8个小时有利于健康和心情舒畅。美国成年人一般每晚睡6～7个小时。一般的久坐不动的上班族并不强调身体的恢复能力。但妈妈们都知道，平均每晚8小时的睡眠有助于恢复，尤其是在非常严格的训练期间。毕竟，睡眠的目的就是调整身体进入恢复状态。睡眠时间越长，恢复质量就越高。

睡眠时间并不完全等同于躺在床上的时间。几乎没有人可以做到头沾枕头就睡着。晚上11点上床、早上7点起床或许无法保证8小时的睡眠。实际一点的做法是，额外增加一些时间来应对入睡的任何延迟，以此保证8小时的睡眠。

\subsection{蛋白质}

运动员到底需要多少蛋白质？最近，越来越多的研究已经阐明了力量训练运动员对蛋白质的需求。美国的每日建议摄入量（Recommended Daily Allowance，简称RDA）认为，15岁以上的男性和女性应当按照0.8克/千克/天的标准摄入蛋白质（即每天每千克体重摄入0.8克蛋白质）。RDA的建议是基于一般人群的，而美国一般人群存在久坐的问题。认为久坐不动的人和正在通过训练计划系统性地增加身体刺激和适应状态的训练者对营养的需求相同是不合逻辑的。事实上，最近的研究表明，虽然久坐人群对蛋白质的需求不高，但他们也没有摄入足够的蛋白质。有详尽的资料证实，任何能够提高肌肉代谢速率的训练同样能够加速肌肉蛋白质的降解和转化速率。研究表明，阻力训练对肌肉蛋白质合成的刺激作用会持续到训练结束后，并且初级训练者的持续时间会比高级运动员更长一些。

肌肉蛋白质合成（MPS），即构建新的肌肉的过程，需要从饮食中摄取蛋白质及碳水化合物（碳水化合物通过协同作用促进这个过程）。有效的训练刺激不可避免地会导致肌肉蛋白质分解，因此肌肉从刺激中恢复并生长，主要是依靠肌肉蛋白质合成的速度超过肌肉蛋白质分解的速度。蛋白质的合成要超过蛋白质的降解，也就是合成代谢或构建过程必须超过分解代谢或降解过程。如果其他蛋白质合成（用来维持或修复损伤组织）所需的营养无法从饮食中充分获取，身体就会从自己的蛋白质储备中获取，而现存的肌肉量就是现成的蛋白质储备。在长期饥饿或慢性刺激状态中，肌肉减少是很正常的过程。身体实际上是在拆东墙补西墙以维持其功能。即使缺乏食源性蛋白质和碳水化合物，训练刺激仍然会以刺激的方式作用于身体。确保从饮食中摄入充足的蛋白质，训练者就获得了为身体合成新的蛋白质所必需的构建模块。如果做不到这一点，新的肌肉蛋白质合成过程就难以发生，这等同于削弱训练效果、浪费宝贵的精力。

那么，到底需要多少蛋白质来支撑训练呢？文献中推荐的参考范围很大，最高可达2.5克/千克/天蛋白质。有些教练和训练者不喜欢精打细算，也不擅长磅和千克之间的换算。有一种简单的，也是被举重和力量训练圈子用了多年的方法，可以确保蛋白质的摄入，即每天每磅体重摄入1克蛋白质，这样一个200磅（90.7千克）体重的运动员应当通过不同的饮食来源每天摄入200克蛋白质。这相当于2.2克/千克/天蛋白质的摄入，这个值虽然超出了普遍推荐的1.2～1.8克/千克/天蛋白质的摄入量，但仍低于文献建议的最高2.5克/千克/天蛋白质的摄入量，同时仍保持了足够高的摄入量，即使少吃一点也能满足完全恢复的需要。这种计算方法没有考虑瘦体重，其前提是身体构成已处于“正常”范围，因此体脂较高的人在计算蛋白质摄入量时应意识到并认真考虑这一点。

这种计算方式同样没有考虑低质量的蛋白质来源和其他热量摄入的影响。如果长期依靠大豆蛋白、大米蛋白、大麻蛋白、豆类蛋白以及其他非动物来源的蛋白质作为主要蛋白质来源的话，就需要摄入比建议量更多的蛋白质，因为这些蛋白质的氨基酸组成不够合理，支链氨基酸（BCAA），即亮氨酸、异亮氨酸和缬氨酸的含量较低。肌肉蛋白质合成也依赖于非蛋白质来源的热量摄入，因此，如果碳水化合物和脂肪的摄入水平较低的话，你需要摄入更多的蛋白质来支持肌肉蛋白质的合成。此外，在高质量的蛋白质食物搭配高质量的碳水化合物和脂肪的情况下，你可以相应降低蛋白质的总摄入量。并且随着年龄的增长，训练者对蛋白质质量会越来越敏感，所以年长的训练者在制订训练计划时需要摄入更多蛋白质，或者摄入更优质的蛋白质，这一点很重要。

很重要的一点是，不管有些保健专家怎么说，没有绝对的证据支持“过量”蛋白质会损伤肾脏的正常排泄功能的说法。实际上，对没有活跃性肾脏疾病的人来说，不存在一个不安全的饮食蛋白质摄入水平。

蛋白质补剂非常有用，能帮助运动员摄取足量的蛋白质，填补饮食蛋白质的摄入量与建议摄入量之间的缺口。并且蛋白质饮料制作简单、饮用方便，有助于训练后的恢复。市面上现有的高质量乳清蛋白补剂，其支链氨基酸的含量和生物利用效率都高于牛肉，当然更是任何豆类蛋白质制品远不能及的。不过，我们吃饭并不只是为了获取支链氨基酸。牛肉以及一份精致的沙拉中还含有高质量乳清蛋白所不具有的其他营养成分。让补剂在训练饮食中居于主导地位听起来可能很诱人，在一些特殊情况下也可能有必要，但最好还是将补剂作为对饮食的补充。如果你需要摄入大量补剂的话，你最好检查一下其他饮食部分的质量如何。

\subsection{热量}

训练期间会消耗热量，这些热量多数来自身体的碳水化合物和脂肪储备，所以，针对训练后恢复的一个很明显的要求就是，增加热量摄入来填补训练的消耗。训练产生热量需求有两个原因：（1）任何形式的、任何训练量和训练强度的训练都会消耗一部分身体的能量储备，这些消耗的热量必须在下一次活动之前予以补充；（2）强度足够的训练会打破体内平衡以及肌肉结构的完整性，从而产生对蛋白质和脂肪/碳水化合物的热量需求，以促进修复和恢复。

训练期间，肌肉倾向于首先使用糖原储备作为燃料，脂肪在阻力训练的过程中起到的供能作用很小。在肌肉的恢复过程中，碳水化合物仍是肌肉蛋白合成的最主要的能量来源。脂肪则成了除肌肉蛋白质合成之外的其他代谢过程的能量来源。在热量和蛋白质摄入足够的情况下，碳水化合物和脂肪的来源并不是非常重要。相比碳水化合物，脂肪需要更长的时间被分解和利用，但训练后身体代谢速率的升高会持续很长时间，在此期间高热量密度（等量的脂肪含有的热量比蛋白质高）的物质在体内缓慢代谢是有益的。除了充足的热量和蛋白质的摄入，饮食构成中还必须考虑维生素、必需脂肪酸和膳食纤维的需求。饮食质量必须在训练者条件允许的范围内保持尽可能高的水平。

训练日的热量总摄入大于热量总消耗非常重要。摄入和消耗持平在理论上可以维持体内平衡和既有力量水准，但无法支持最大限度的力量增长和肌肉量增长，而这才是力量训练的目的。从实际出发，取得进步需要明显的热量盈余，但每日消耗的热量几乎不可能精确地计算出来，因为变量太多了。训练负荷、睡眠、训练者的性别、饮食本身、训练者的年龄和生长状态等都会影响热量的消耗。如果只是补充日常活动和训练所消耗的热量，我们就无法为体内平衡的恢复和肌肉蛋白质合成产生的适应状态提供所需的额外能量。为了变得更强壮，传统文献建议额外摄入200～400千卡的热量。这对严肃对待力量训练的人来说严重不足，对体重过轻想要增加肌肉量的人来说更是远远不够。作者本人的经验是，每天能够保证身体恢复和力量增长的较为合理的热量摄入至少要比基准需求高出1000千卡。如果主要目标是增肌的话，你至少要保证每日的热量摄入比基准需求高出2000千卡，对代谢效率较低的人来说，数值还要更高一些。考虑到计算基础消耗基本不可能实现，最好的实践方法就是进食尽可能多的植物和动物来源的高质量蛋白质、碳水化合物和脂肪，然后根据摄入结果进行调整。对想要增重的年轻运动员来说，全脂牛奶永远是优质的补充。

\subsection{脂肪酸}

必需脂肪酸（Essential Fatty Acids，EFAs）是另一种影响身体恢复的物质。虽然有些圈子对膳食脂肪的偏见仍很流行，但脂肪是必需的营养物质，也是高效的热量来源。不存在必需碳水化合物，但存在几种必需脂肪酸——尽管身体可以以食物中的脂肪为原料合成多种脂类，但无法合成$\omega$-3和$\omega$-6脂肪酸这样的多不饱和脂肪酸。这两种脂肪酸对身体结构的完整性起着极为重要的作用，对免疫功能和视敏度至关重要，并且参与类花生酸（Eicosanoids）——能够调节炎症过程的前列腺素（Prostaglandins）的前体——的合成。其中，$\omega$-3脂肪酸对身体恢复最为重要：它能够支持合成代谢过程，辅助控制训练后的炎症和疼痛。并且$\omega$-3脂肪酸在饮食中也较为匮乏。错误比例的$\omega$-6脂肪酸摄入反而会加剧炎症过程。

必需脂肪酸摄入不足的情况在美国较为普遍，因为$\omega$-3脂肪酸的主要来源——鱼类，从未成为美式饮食的重要组成部分。必需脂肪酸的长期严重匮乏会导致身体发育迟缓、皮肤干燥、腹泻、伤口愈合缓慢、感染率升高和贫血。亚临床匮乏不会产生通过观察就能得出诊断的症状。执行极端低脂饮食的人可能会突发急性临床缺乏症，症状在2～3周内就会明显显现。

其实，只要几克富含$\omega$-3脂肪酸的油就够了。每天饱餐一顿富含脂肪的鱼，比如三文鱼，就可以解决问题。很多人发现，每天吃一些含有$\omega$-3脂肪酸的鱼油补剂很有帮助，因为较高水平的必需脂肪酸摄入量对严苛的训练大有裨益。鳕鱼肝油是一种廉价的必需脂肪酸来源，同时富含维生素A和维生素D。

\subsection{水分}

为了从繁重的训练中恢复，水分是必不可少的。毕竟，人体内几乎所有的生物化学过程都发生在水介质中。缺水会导致运动表现下滑，严重的话，会带来灾难性的后果。随着代谢率的升高，对水的需求也会增加。增加肌肉内的能量底物的储存（比如ATP、磷酸肌酸和糖原）也会增加对细胞内水分的需求。水分充足的细胞才能进行合成代谢。事实上，一个细胞，尤其是多核肌肉细胞，其在缺水状态下合成蛋白质的速率会远低于充水状态下的速率。肌酸补剂促进骨骼肌增长的方式之一就是提高细胞含水量。那么，我们需要喝多少水才能帮助恢复、避免过度训练呢？

每个人身边都有医生、营养师、饮食学家、训练员、教练或者朋友“知道”你“绝对应该”按照“8×8”的方式喝水：每天8杯水，每杯8盎司（236.6毫升）。这相当于每天喝0.5加仑或1.9升水。注意，标准的饮料易拉罐是12盎司（354.8毫升）或20盎司（591.4毫升）的，餐馆里的“小杯”一般是16盎司（473.1毫升），所以这个建议并不是指常见的8“杯”。

但我们真的需要喝这么多水吗？“8×8”的建议并不是基于科学研究结论得出的，而是1974年由临床专业人员撰写的营养文本提出的主观看法，之后逐渐变成了牢不可破的临床教条和“传统智慧”。大部分研究数据指示的液体摄入量为每天1.2～1.6升（比8×8法中的1.9升要少），这已足够维持一个健康、温和锻炼的人对水的需求了。这些建议当然需要考虑不同的环境状况——6月在佛罗里达州的人的饮水需求与10月在加拿大马尼托巴省（Manitoba）的人显然是不同的——所以在现实生活中，并没有绝对的摄入参考量。自从瓶装水工业诞生以来，很难想象谁的身体无法自发地调节其体内的水合状态。毕竟，人类社会很少会出现每天背着或拿着水时不时喝几口的习俗。因此，大多数时候，感觉口渴就自己补充点水可能是一种合适的维持身体健康和机体功能的方法。但这样的摄入量足以支持从高强度训练中恢复的需要吗？

每天1.2～1.6升水据说能够支持温和运动的生活方式，那么在这个基础上训练者还需要喝多少水呢？答案是：要多得多。活跃训练的人需要喝更多的水，才足以支持他们数量增加的代谢活跃组织、逐渐增加的训练负荷导致的更多热量消耗以及相对低效的散热系统。在饮水方面，没有一个统一的标准。一个很好的经验法则是，每消耗1000千卡热量就喝1升水。一天消耗5000千卡的话就需要喝5升水（或1.3加仑）。这个量看起来可能很大，但考虑到一个非常活跃的运动员每天训练数小时的需求，这个建议量还是合理的，并且如果在温度较高的环境中训练的话，这个量或许还不够。每天超过5升的饮水相当多了，所以你需要关注水的摄入。\textbf{过量饮水极其危险，由此导致的低钠血症（Hyponatremia）有潜在的致命风险，}但喝下可以达到毒性剂量的水需要刻意尝试，需要付出远远超过只为缓解口渴时的努力，一般人也不可能无意地喝下这么多水。然而，这种情况在一些民间组织的耐力赛事中偶有发生，由于活动中的每个休息站点都有饮水提供，过分热情却毫无经验的参赛者错误地听从了所谓的“专业人士”未经深思的建议，为了永远比脱水“更快一步”，在口渴之前就开始喝水。

最后一点需要考虑的是，使用何种液体形式完成补充。很多健康从业者会大胆声称，只有水和少数其他的“自然”饮料才算数。任何含有咖啡因、酒精甚至糖的饮料都不被他们当作合适的补水饮品。他们的原话是“你不会用软饮料来洗车，那为什么要用它们来补水？”这种荒谬的言论显示出他们对水分吸收机制缺乏了解。水分是从肠道内被吸收的，任何含水的饮品，甚至含水量高的食物，如果其提供的水分比用来代谢它的水分要多的话，都算作水分的摄入。水本身是最好的补水液体，毕竟，水就是水。如果基于正常补水条件下的实际考虑，水的吸收要比常见的商业饮料更快，但每种基于水的饮料（难道还有不是基于水的饮料？）都能帮助补水。一罐20盎司（591.4毫升）的健怡可乐也能补水，虽然其中含有咖啡因和人工甜味剂；一罐含有高果糖玉米糖浆的普通可乐也能补水，即使里面含有咖啡因和糖。酒精饮料在人类历史上一直是一种非常高效的水合剂。在比较早的时代，啤酒和红酒是维持人们存活下去的主要补液，因为它们比当时未经净化的水更为安全，就像18世纪英国海军的格罗格酒（Grog）——由1份朗姆酒和4份水混合，再加一点青柠汁制成。我们并不是建议把软饮料、啤酒和红酒当作训练饮食中的固定组成，只是诚实地考虑美国生活的现实状况以及其对恢复的影响。抛开这些饮料的其他优缺点不谈，适当饮用的确可以补充水分。

\subsection{维生素和矿物质}

经常有人说，普通的美式饮食能够提供健康生活所需要的全部维生素和矿物质，这种认识也总是把努力训练的人群包括进去。除非已经显示出了缺乏症状，否则几乎没有人会去测试机体的维生素和矿物质的水平。因此，不论是习惯久坐的人还是经常运动的人，几乎没有人能够确定他们摄入了足够的维生素和矿物质。

在美国，严重的维生素和矿物质缺乏的情况并不常见，但也不是不存在。轻度缺乏的情况则常见得多，比如，大部分美国女性存在长期缺铁和缺钙的情况，尽管程度并不严重，但症状确实存在。钙在神经和肌肉的生理功能、生长过程和运动表现中发挥多种重要作用，缺钙会使训练恢复受到限制。铁在氧运输和代谢功能中扮演重要角色，轻度缺铁会对身体训练后的恢复能力产生明显的负面影响。

维生素和矿物质在体内的生化反应中发挥中介物的作用。它们被称为“微量营养素”，身体对其需求量的确很小，它们在食物中的自然含量也不同。为生活计，当然也是为了训练，为了获得所有的维生素和矿物质，我们需要食用多种食物。一般的美国孩子做不到这一点。如果父母有意识地为年轻运动员提供高质量、种类丰富的饮食的话，这或许就不是个问题。但美国的文化特色是便利和嗜好，通常人们只有少数一些食物供选择，并且为了便于储存和烹饪，它们还都经过了加工。这个加工过程一般会在某种程度上减少食物中的维生素和矿物质含量，会使饮食的质量变差，即使饮食中含有充足的热量和蛋白质。

结果是，尽管典型的运动员饮食中的维生素和矿物质含量不至于不足到出现症状的地步，但很可能达不到从严酷的训练中恢复身体所需的最佳含量。最近的研究表明，维生素补剂对延长习惯久坐的人的寿命没有任何作用，当然我们也不关心长寿，相比之下，我们更关心运动表现。毕竟，如果一个刻苦训练的运动员的热量、水分、蛋白质需求都提高了的话，那么摄入更多的维生素和矿物质也是理所当然的。美国人口的地域性、宗教、文化以及经济品位和习惯都很多样化，针对性的测试又非常昂贵，运动员或教练没有简单易行的方法来评估饮食中的维生素和矿物质含量。谢天谢地，也并不需要这样：营养补剂可以确保这些训练和恢复所需的要素得到安全、有效和经济的补充。

最好的方法既简单又便宜。包含所有常见维生素和矿物质的廉价的通用型复合片在杂货店和网上都有销售。多花些钱能让你买到品质更好、纯度更高、吸收更容易的产品。比尔·斯塔尔（Bill Starr）在他著名的《强者生存》（\textit{The stongest shall survive}）一书中提出了“铲子方法”：吃很多，身体会把不需要的部分排出来。维生素中毒是极其罕见的，尤其是对高强度训练的运动员来说，所以这个方法值得一试。

\section{训练强度和训练量}

\subsection{周期化}

50年前，或许是受到了塞利理论的启发，人们意识到可以将该理论直接应用于运动员的训练，前苏联的运动生理学家因此提出了几种可以充分利用身体适应渐增负荷能力的训练方法。这种方法的根基叫作\textbf{周期化}（Periodization），常被归功于20世纪60年代、前苏联时期的列昂尼德·马特维耶夫（Leonid Matveyev）（苏式周期化的高级版本最早可以追溯到20世纪40年代和50年代的匈牙利）。到了20世纪70年代，卡尔·米勒（Carl Miller）将周期化训练的概念引入了美国举重界；1981年，迈克·斯通（Mike Stone）以此为基础提出了一个提高运动表现的力量训练模型。从那时开始，不管什么运动项目，周期化成为成功制订训练计划的主要工具。

所有提高运动表现的训练的目的都应是使身体通过塞利刺激模型中的第1阶段和第2阶段，提供足够的训练刺激来引发身体的适应状态，同时避免第3阶段——力竭的发生。正确设计的训练计划会通过调节训练量和训练强度来控制身体承受的刺激，从而得到良好的结果。因此，量化训练量和训练强度非常重要。

\textbf{训练量}是训练者在一次或一组训练中举起的总重量，即

$$\text{重复次数}\times \text{单次举起的重量} = \text{训练量}$$

下表是一个计算深蹲训练训练量的简单示例。

\begin{table}[htbp]
    \centering
    % 定义9列：增加第4列和第8列作为空列（c），形成微小的自然间隔
    \begin{tabular}{rrr c rrr c l}
        \multicolumn{3}{l}{热身组} & & \multicolumn{3}{l}{正式组} & & \\
        \cline{1-3} \cline{5-7} % 横线会自然在第4列处断开
        45 & 95 & 135 & & 185 & 185 & 185 & & 重量 \\
        5 & 5 & 5 & & 5 & 5 & 5 & & 重复次数 \\
        225 & 475 & 675 & & 925 & 925 & 925 & & 每组的训练量 \\
        \cline{1-3} \cline{5-7} % 横线会自然在第4列处断开
        \multicolumn{3}{r}{1375} & & & & & & 热身组的训练量 \\
        & & & & \multicolumn{3}{r}{2775} & & 正式组的训练量 \\
        \cline{5-7} % 仅在5-7列下方画横线
        & & & & \multicolumn{3}{r}{4150} & & 总训练量 \\
    \end{tabular}
\end{table}

在这次训练中，包括热身组在内，训练者深蹲的总重量是4150磅（1882.4千克）。训练中每次动作重复都被计算在内，这样作为所施加的刺激的总训练量就可以被量化了。通常，只考虑正式组会更有用一些，因为打破体内平衡、引发第1阶段的是正式组，而非热身组。正如上面的表格所示，这样算下来训练量减少了很多。但如果训练者的热身组数非常多，则需要认真考虑其对总训练量的影响。

\textbf{强度}是在一次或一组训练中使用的平均重量相对于训练者1RM的百分比。

\begin{align*}
    \frac{\text{训练量}}{\text{重复次数}} &= \text{使用的平均质量}\\
    \frac{\text{使用的平均重量}}{1\text{RM}}\times 100\% &= \text{训练强度的百分数}\%
\end{align*}

在上面的例子中，使用的平均重量是4150磅（1882.4千克）/30次，即138.3磅（62.7千克）/次。如果训练者的1RM是225磅（102.1千克），那么其训练强度就是138.3/225×100\%=61\%。这就很容易看出来，将热身组计算在内对平均每次使用的重量和平均强度有多大的影响了，因此只用正式组计算训练强度就可以了。如果上面的例子只计算正式组的话，训练强度为82\%。

我们重新叙述一遍：训练强度是相对于1RM的百分比。1RM的80\%的训练强度大于1RM的50\%。这个概念很简单，但在科学、医学和流行文献中对于“训练强度”的定义有多种方式。有时候，训练强度等同于某个给定动作中的爆发力输出水平。训练者完成某次特定重复时像专注度一样抽象的事物（比如，“让我们专注下一次重复，提高强度！”）或者对练习的主观努力感受（比如，伯格的“自觉运动强度分级量表”），也会被用来定义强度。还有一种表述与疲劳相关：如果训练使肌肉疲劳，那么训练就是有强度的。这些概念都在文献中出现过，可能比较适合耐力训练这种训练刺激本质上依靠较小训练刺激产生的累积效应的运动。但对职业力量训练来说，这些定义都无一例外地不切实际，因为这些概念无法量化，而科学家和参与者都将量化特征看作关键要素。把强度定义为相对于1RM的百分比，这种做法或许看起来过于简单，但这正是其优势所在。

这是最实际、最有用的方法，对要给一大群人制订训练计划，需要一种客观地评估训练强度和训练进步方法的教练和训练员来说尤其如此。

训练重量相对于1RM的训练强度范围的简单计算如下。

\begin{table}[htbp]
    \centering
    \begin{tabular}{l r c l}
        \hline
        \multicolumn{2}{l}{深蹲强度 1RM（磅）} & & 225 \\
        \hline
        95\% & $225 \times 0.95$ & $=$ & 214 \\
        90\% & $225 \times 0.90$ & $=$ & 203 \\
        85\% & $225 \times 0.85$ & $=$ & 191 \\
        80\% & $225 \times 0.80$ & $=$ & 180 \\
        75\% & $225 \times 0.75$ & $=$ & 169 \\
        70\% & $225 \times 0.70$ & $=$ & 158 \\
        \hline
    \end{tabular}
\end{table}

通常，根据训练者的不同水平，可以将训练划分为时间长度和负荷特征不同的几个阶段，以此来实现对周期化训练中的训练量和训练强度——身体承受的刺激程度——的控制。

解读有关过度训练的文献时，必须清楚，其中很多研究是基于有氧训练的。像举重这样的无氧训练导致的过度训练与有氧训练导致的过度训练有着很大差别，因为两者的刺激类型完全不同——力量训练者将有氧训练称为慢速长距离训练（Long Slow Distance，简称LSD）。有氧和无氧训练对训练量和训练强度的不同定义方式会影响对过度训练的分析。比如，现代公路自行车选手可能每天都会在自行车训练上花费数小时，由此产生的训练量是巨大的。如果他们想练得更努力，他们可以增加里程、每天的训练时间或者训练天数，累积所谓的“垃圾里程”。他们也会按照最大摄氧量（VO$_{2}$max）的某个可持续的百分比来骑行，如果用这种方法来衡量公路骑行运动的强度，那么每次训练的平均强度都会与先前差不多。需要指出的是，最大摄氧量会出现在肌肉达到30\%～40\%最大收缩力量的时候。因此，强度（量化为绝对力量的百分比）在一般的自主训练的美国竞技自行车手的训练计划中并不是主要因素。因为调节的训练变量主要是训练量，公路自行车手往往更容易出现基于训练量的过度训练。但你必须清醒，自行车手的“训练量”与举重选手的“训练量”是两个完全不同的概念。按照我们的定义方式，骑行的训练量通常是在相对低得多的强度下进行的，但其涉及的重复次数则要比任何杠铃训练高上千倍。

因为力量训练计划需要兼顾对训练量和训练强度的控制，所以训练者有机会体验分别由两者引起的过度训练。极端的训练风格，比如最具代表性的“要么放弃，要么选择大重量”的方法，会导致基于训练强度的过度训练，而另一个极端，“练至力竭”法，可能会导致基于训练量的过度训练。因为大多数的训练计划都会涉及这两种变量，所以由二者共同导致的过度训练更为常见。

理解阻力训练中两种刺激下的恢复速率极为重要。由训练量引发的过度训练主要影响肌肉细胞的收缩组件和代谢系统，而由训练强度引发的过度训练主要影响神经系统的功能以及其与肌肉系统的交互作用。对训练有素的举重运动员来说，由训练强度引发的过度训练更容易恢复。当为力量型或者爆发力型比赛做准备时，随着训练强度的持续增加，训练量则会大幅削减，直到比赛临近。当为耐力型比赛做准备时，力量训练的训练量和训练强度都应提前几周进行削减，因为基于训练量的恢复需要较长的时间，会直接影响比赛成绩。

针对过度训练的最基本的处方，是让运动员投入时间恢复，同时削减训练负荷。处理过度训练的时间影响了教练和训练者的宝贵进步：削减负荷不会产生任何进步，甚至无法维持原有水平；完全停练则会不可避免地导致训练不足。既然过度训练的代价这么大，最好的方法还是防患于未然。正确制订适合运动员和运动项目的训练计划是关键。尽管执行妥当的简单线性训练计划能使运动员在早期取得极快的进步而不会引起过度训练，但是对水平更高的运动员来说，更为复杂的训练计划——周期化训练计划——是必不可少的。